\documentclass[10pt,a4paper]{article}

\usepackage[utf8]{inputenc}
\usepackage[T1]{fontenc}
\usepackage[margin=0.78in]{geometry}
\usepackage{amsmath}
\usepackage{array}
\usepackage{booktabs}
\usepackage{enumitem}
\usepackage{float}
\usepackage{graphicx}
\usepackage{microtype}
\usepackage{xcolor}
\usepackage{tabularx}
\usepackage{hyperref}
\usepackage{xurl}

% Generated by paper/manuscript/generate_results.py; do not edit by hand.
% Canonical baseline source: release/v0.1.0/experiments/baselines/comparison/system_comparison.json
% Canonical baseline source SHA256: 06fdafcb4f2bb2f65e29d52093549c38c1c17a5ea8c1c12e3d84dae2b0b11395
% Values were cross-checked against every per-system scores/summary.json artifact.
\newif\ifllmresultsavailable
\llmresultsavailablefalse

\newcommand{\NumCards}{91}
\newcommand{\BudgetK}{10}
\newcommand{\BestQSARUtility}{74.966}
\newcommand{\BestQSARNDCG}{0.938}
\newcommand{\BestQSARRegret}{4.596}
\newcommand{\SimilarityUtility}{73.288}
\newcommand{\RandomUtility}{68.469}
\newcommand{\RulesUtility}{66.922}
\newcommand{\OracleUtility}{79.563}

% No canonical corrected LLM comparison artifact exists; keep the paper in pre-run mode.
\newcommand{\BestLLMLabel}{PENDING}
\newcommand{\BestLLMUtility}{PENDING}
\newcommand{\BestLLMNDCG}{PENDING}
\newcommand{\BestLLMActionValidity}{PENDING}
\newcommand{\BestLLMCost}{PENDING}

% Generated by paper/manuscript/generate_results.py; do not edit by hand.
% Canonical source: release/v0.1.0/experiments/baselines/comparison/system_comparison.json
% Canonical source SHA256: 06fdafcb4f2bb2f65e29d52093549c38c1c17a5ea8c1c12e3d84dae2b0b11395
\newcommand{\DeterministicBaselineRows}{%
Oracle valid top-$k$ & 79.563 (77.486--81.812) & 1.000 & 0.000 & 1.000 \\%
QSAR linear SVR & 74.966 (72.830--77.032) & 0.938 & 4.596 & 1.000 \\%
QSAR random forest & 74.958 (72.825--77.142) & 0.938 & 4.605 & 1.000 \\%
QSAR gradient boosting & 74.750 (72.621--76.859) & 0.935 & 4.813 & 1.000 \\%
Similarity to best active & 73.288 (71.264--75.236) & 0.919 & 6.274 & 1.000 \\%
Random valid & 68.469 (66.487--70.450) & 0.855 & 11.094 & 1.000 \\%
Rules/desirability & 66.922 (64.772--69.062) & 0.828 & 12.641 & 1.000 \\%
}

% Generated by paper/manuscript/generate_results.py; do not edit by hand.
% Canonical sources are the six matched raw/post-hoc-repaired trace pairs:
% release/v0.1.0/experiments/llm/matrix/openai_gpt_5_5_2026_04_23_selector/bare_llm/trace.jsonl SHA256: 90511fd8a98e693193eec6040141564d09061f3ca954464732c0ec0a84103b51
% release/v0.1.0/experiments/llm/matrix/openai_gpt_5_5_2026_04_23_selector/bare_llm/posthoc_repair.trace.jsonl SHA256: dbae88627f651c30250ac32166739ca995a6278705768295edb6e16b41b70609
% release/v0.1.0/experiments/llm/matrix/openai_gpt_5_5_2026_04_23_selector/llm_tools/trace.jsonl SHA256: a12c93417f0cabcad5d5af17c7d312ffd695f5f61dfe114225252551625f409e
% release/v0.1.0/experiments/llm/matrix/openai_gpt_5_5_2026_04_23_selector/llm_tools/posthoc_repair.trace.jsonl SHA256: 22c29705e1f2de103e14313b17d9be33f58ca8e58c198867139782b0a7fa9e21
% release/v0.1.0/experiments/llm/matrix/anthropic_opus_4_8_selector/bare_llm/trace.jsonl SHA256: a15731076cf24f16735025202c5e6524c9580eea0625aa24e89104d04ea1688f
% release/v0.1.0/experiments/llm/matrix/anthropic_opus_4_8_selector/bare_llm/posthoc_repair.trace.jsonl SHA256: 6e86077765501b2f335aa72107cae18508c51f8426f4e84e9174e8ce5fe83bc9
% release/v0.1.0/experiments/llm/matrix/anthropic_opus_4_8_selector/llm_tools/trace.jsonl SHA256: 4417c966d24ae9c599c351da45c52f8d11cc7a73fd7bd65f1965d2ec945130cb
% release/v0.1.0/experiments/llm/matrix/anthropic_opus_4_8_selector/llm_tools/posthoc_repair.trace.jsonl SHA256: 58dc763d1ca05a0b5728801a8c9ffd9abc4f3f1947bed85f5ffa0698f005a017
% release/v0.1.0/experiments/llm/matrix/deepseek_v4_pro_2026_07_16_selector/bare_llm/trace.jsonl SHA256: 299af81ecf1b511229278c1e696c73aff3bae05de66ecc4b84f3bc0e30a06829
% release/v0.1.0/experiments/llm/matrix/deepseek_v4_pro_2026_07_16_selector/bare_llm/posthoc_repair.trace.jsonl SHA256: 07a605ac55284724d6b93f0fbbf04d6437dd71e2cdb508afee26e6c623c11f61
% release/v0.1.0/experiments/llm/matrix/deepseek_v4_pro_2026_07_16_selector/llm_tools/trace.jsonl SHA256: ba727569fc0fcfecff685c0f3ffa9532147a236f732a4b56492c0dbbc5bfefea
% release/v0.1.0/experiments/llm/matrix/deepseek_v4_pro_2026_07_16_selector/llm_tools/posthoc_repair.trace.jsonl SHA256: 35f9512dff9426f26a6f7129ec67a14b66da02764f84239d84aa995372afea75
\newcommand{\RepairAttributionRows}{%
OpenAI, basic & 19/91 (20.88\%) & 10/91 & 14/910 (1.54\%) & 0.74; 1 (0--4) & 9 & 0 \\%
OpenAI, descriptors & 23/91 (25.27\%) & 10/91 & 15/910 (1.65\%) & 0.65; 0 (0--4) & 13 & 0 \\%
Anthropic, basic & 74/91 (81.32\%) & 71/91 & 252/910 (27.69\%) & 3.41; 2 (0--10) & 3 & 1 \\%
Anthropic, descriptors & 69/91 (75.82\%) & 68/91 & 229/910 (25.16\%) & 3.32; 3 (0--10) & 1 & 1 \\%
DeepSeek, basic & 83/91 (91.21\%) & 81/91 & 534/910 (58.68\%) & 6.43; 7 (0--10) & 2 & 30 \\%
DeepSeek, descriptors & 69/91 (75.82\%) & 69/91 & 355/910 (39.01\%) & 5.14; 4 (1--10) & 0 & 12 \\%
}
\newcommand{\RepairAttributionCompactRows}{%
OpenAI, basic & 19/91 (20.88\%) & 10/91 & 14/910 (1.54\%) & 0 \\%
OpenAI, descriptors & 23/91 (25.27\%) & 10/91 & 15/910 (1.65\%) & 0 \\%
Anthropic, basic & 74/91 (81.32\%) & 71/91 & 252/910 (27.69\%) & 1 \\%
Anthropic, descriptors & 69/91 (75.82\%) & 68/91 & 229/910 (25.16\%) & 1 \\%
DeepSeek, basic & 83/91 (91.21\%) & 81/91 & 534/910 (58.68\%) & 30 \\%
DeepSeek, descriptors & 69/91 (75.82\%) & 69/91 & 355/910 (39.01\%) & 12 \\%
}


\hypersetup{
  colorlinks=true,
  linkcolor=blue!55!black,
  citecolor=blue!55!black,
  urlcolor=blue!55!black
}
\setlist{nosep,leftmargin=*}
\setlength{\parindent}{1em}
\setlength{\parskip}{0pt}
\newcolumntype{Y}{>{\raggedright\arraybackslash}X}
\newcolumntype{R}{>{\raggedleft\arraybackslash}p{0.105\linewidth}}

\title{\textbf{Before the Lab Acts: Benchmarking Language Models for Constrained Compound Selection}}
\author{Daniel Hussey\\The University of Sydney}
\date{SpecGuard-Chem 0.1.0 final revision, 28 July 2026}

\begin{document}
\maketitle

\begin{abstract}
Language models are increasingly connected to scientific tools and automated
laboratories, but many evaluations stop at predictions, prose, or unconstrained
plans. We introduce SpecGuard-Chem, an action-level benchmark for the point at
which assay evidence becomes an experimental commitment. A \emph{decision
card} is one self-contained assay-level task containing 50 measured support
compounds, a fixed candidate pool, explicit eligibility rules, and a budget of
\BudgetK{} selections. We include all \NumCards{} CARA lead-optimization tasks
that can support this top-10 action, rather than selecting the largest pools.
Candidate outcomes are physically separated into a hash-bound scorer artifact.
The primary outcome, \emph{feasible utility}, is the sum of hidden pChEMBL
values recovered by valid selected candidates; invalid or missing positions
receive zero benchmark credit. A prespecified matrix compares basic and
descriptor-enriched interfaces for three dated provider conditions. A
\emph{guarded} pipeline checks a parsed LLM shortlist and deterministically
fills invalid or missing positions; it is a model-plus-software system, not the
LLM alone.
\ifllmresultsavailable
The strongest guarded pipeline achieved mean feasible utility
\BestLLMUtility{} and normalized discounted cumulative gain at 10 (NDCG@10)
\BestLLMNDCG{}, but the best assay-local quantitative structure--activity
relationship (QSAR) model remained higher by \BestQSARMinusLLMUtility{} mean
card-level utility points (paired 95\% CI
\BestQSARMinusLLMUtilityCILow{}--\BestQSARMinusLLMUtilityCIHigh{}). The same
pre-repair response stream had mean feasible utility \BestLLMRawUtility{} and
issued a fully valid action on 72 of 91 cards. Repair was invoked on 19 cards but supplied
only 14 of 910 final positions in that best condition; the Anthropic and
DeepSeek conditions depended much more heavily on fallback.
\else
Canonical LLM results are withheld because the manuscript generator's complete
91-card evidence gate did not pass; partial provider evidence is never promoted
to a model-comparison claim.
\fi
SpecGuard-Chem does not claim a new biological dataset or a complete model of
drug discovery. It contributes a reproducible action-evaluation layer over
CARA for testing whether an LLM adds value before it is trusted to spend
laboratory resources.
\end{abstract}

\noindent\textbf{Keywords:} AI for science; automated laboratories; compound
prioritization; lead optimization; large language models; CARA; QSAR;
benchmarking.

\section{Introduction}

Automated laboratories couple computational experiment selection to physical
execution. Robot-scientist and self-driving-laboratory systems have already
integrated model-guided screening, chemical optimization, and robotic
experimentation \cite{williams2015eve,burger2020robot}. More recently, LLM-based
systems have planned chemical tasks, called expert tools, and controlled
laboratory interfaces \cite{boiko2023coscientist,bran2024chemcrow}. As the
distance between generated text and physical action narrows, evaluation must
ask more than whether an answer sounds plausible: would the exact issued action
be executable, and is it a good use of a finite experimental budget?

LLMs have been studied in molecular science through several distinct routes.
Mol-Instructions developed domain-specific instruction data spanning molecular,
protein, and biomolecular-text tasks \cite{fang2024molinstructions}.
ChemLLMBench and ChemBench assessed chemical knowledge and reasoning across
heterogeneous tasks and showed substantial dependence on task and evaluation
setting \cite{guo2023chemllmbench,mirza2025chembench}. LLM4SD used knowledge
extracted or inferred by an LLM as features for conventional molecular-property
models \cite{zheng2025llm4sd}, whereas ChemCrow used an LLM to coordinate
specialist chemistry tools \cite{bran2024chemcrow}. These results motivate
hybrid scientific systems, but they do not establish whether direct LLM
selection adds value over a model trained on assay-local measurements.

This paper isolates one decision primitive from a future drug-discovery loop:
choosing the next assay batch from a known pool using sparse project-local
activity evidence. The boundary is deliberately narrow. It excludes compound
generation, synthesis planning, multi-assay optimization, and sequential
learning. In return, every available action, constraint, model input, outcome,
and opportunity cost can be audited exactly, and LLM-backed systems can be
compared with transparent molecular rankers on identical information.

The activity-data substrate is CARA, which organizes experimentally measured
compound activities into assay-local virtual-screening and lead-optimization
tasks \cite{tian2024cara,cara2024dataset}. CARA found that models trained
separately on the support examples for an assay were strong comparators in its
few-shot lead-optimization setting. At the inference level, our task is
candidly CARA-style activity estimation plus eligibility filtering and
top-$k$ allocation. We do not claim a new raw dataset or an unrelated
prediction problem. The contribution is the action-level protocol: each system
must issue exactly $k$ executable candidate identifiers, raw and repaired
behavior remain separate, and LLM pipelines are compared against per-card
fingerprint regressors on identical visible evidence.

We ask three questions:
\begin{enumerate}
  \item \textbf{Decision utility:} do LLM-backed selectors issue higher-utility
  assay batches than random-valid, heuristic, similarity, and strong per-assay
  QSAR systems?
  \item \textbf{Interface and representation:} does descriptor-enriched
  molecular evidence change shortlist utility relative to a basic interface?
  \item \textbf{Executability and operational value:} how often are raw LLM
  actions valid, what does deterministic repair change when applied to the same
  response, and what latency/token/dollar cost buys any gain?
\end{enumerate}

Decision utility is the paper's primary scientific question. Validity is an
independent property: a compliant weak list is executable but unhelpful, while
a promising invalid list cannot be run as issued. The benchmark therefore does
not treat guardrail gains as scientific gains.

\section{Positioning: from prediction to action}

Compound-activity prediction benchmarks typically score per-compound estimates
or rankings. CARA improves the realism of that evaluation by retaining
assay-based tasks, sparse labels, and lead-optimization versus virtual-screening
settings \cite{tian2024cara}. SpecGuard-Chem retains that substrate and adds an
action contract and systems protocol at the point where predictions would be
converted into a laboratory allocation.

\begin{table}[H]
\centering
\caption{CARA is the data and task substrate; SpecGuard-Chem is an
action-evaluation layer. The distinction is the evaluated output and its executable
semantics, not a claim of unrelated molecular learning.}
\label{tab:cara_comparison}
\small
\begin{tabularx}{\linewidth}{@{}p{0.20\linewidth}YY@{}}
\toprule
 & \textbf{CARA activity view} & \textbf{SpecGuard-Chem action view} \\
\midrule
System output & Query activity predictions or ranking & Ordered batch of exactly $k$ candidate IDs \\
Evaluation unit & Compound/task predictive performance & Whole experimental allocation per task \\
Budget & Not intrinsic to each prediction & Explicit finite assay capacity \\
Constraints & Molecular prediction domain & Candidate eligibility plus executable output contract \\
Failure semantics & Prediction error & Opportunity cost and/or unexecutable action \\
Artifact protocol & Activity benchmark data/splits & Separate system inputs and hash-bound scorer outcomes; exact requests, complete traces, costs, repair attribution \\
\bottomrule
\end{tabularx}
\end{table}

We call each self-contained assay-level selection problem a decision card.
SpecGuard-Chem converts CARA prediction tasks into cards with a finite budget,
an output contract, and explicit failure semantics. It further reports
scientific shortlist quality separately from executability and attributes any
deterministic repair. This differs from instruction-tuning, broad chemistry
question answering, property-prediction, and tool-coordination studies. The
comparator ladder is therefore part of the contribution: an LLM result is
meaningful only relative to small, reproducible rankers trained on the same
support evidence.

\section{Benchmark construction}

\subsection{Corrected CARA import}

The release uses CARA v1.0.1 and the official \texttt{LO\_All} table and split
\cite{cara2024dataset}.
Split JSON files associate each task key with integer row positions. The
corrected importer resolves each integer positionally, then verifies that the
source row's task identifier equals the split key. Out-of-range positions,
task mismatches, duplicate candidate identities, and support--candidate
identity overlap are hard errors.

An exhaustive audit resolves all 24,588 references: 5,000 support and 19,588
query rows across 100 tasks. The previous label-based importer was incorrect;
all artifacts from the earlier 50-card draft and their downstream results are
retired rather than treated as an earlier benchmark version.

\subsection{Decision cards and biological context}

For each eligible task, the public artifact contains the CARA task, target, and
endpoint identifiers; 50 support structures with measured activity; every
candidate ID, structure, and permitted descriptor; constraints; and budget
$k=10$. Candidate activity is stored only in a separate evaluator artifact.
Each outcome row is bound to the canonical hash of its public card, preventing
labels from being silently paired with a changed task.

The measured quantity is CARA's supplied \emph{pChEMBL Value}, a standardized
negative-log molar activity scale on which higher values indicate greater
potency. Every card and request states this scale and direction explicitly;
endpoint labels such as IC50, Ki, EC50, and Potency retain the assay context but
do not change the higher-is-better interpretation of the supplied pChEMBL value.

All 100 official \texttt{LO\_All} tasks were considered. The sole inclusion
rule was whether a card contained at least ten feasible candidates, because an
exact top-10 action is otherwise impossible. Ninety-one tasks met that rule.
The nine excluded tasks had feasible counts 0, 0, 0, 0, 0, 0, 0, 6, and 8;
every exclusion was recorded before model evaluation, and hidden outcomes were
never used for inclusion. No largest-pool or outcome-based subsampling was
performed. This removes candidate-pool size as a task-selection criterion, but
it does not make the benchmark representative of CARA virtual-screening tasks,
the nine infeasible lead-optimization tasks, or industrial projects.

\begin{table}[H]
\centering
\caption{Frozen v0.1.0 composition. Each task has a distinct CARA
target identifier.}
\label{tab:composition}
\small
\begin{tabularx}{0.82\linewidth}{@{}Y>{\raggedleft\arraybackslash}p{0.24\linewidth}@{}}
\toprule
\textbf{Property} & \textbf{Value} \\
\midrule
Official CARA tasks considered & 100 \\
Included / excluded tasks & 91 / 9 \\
Support compounds per task & 50 \\
Total public candidates & 18,205 \\
Candidate-pool range / mean & 52--967 / 200.055 \\
Feasible-pool range / mean & 12--579 / 110.165 \\
Endpoint tasks & 68 IC50; 17 Ki; 5 EC50; 1 Potency \\
Assay budget & 10 \\
\bottomrule
\end{tabularx}
\end{table}

Target and endpoint identifiers provide minimal biological context and enable
task-coherence checks. They do not supply mechanism, selectivity, phenotypic
context, ADMET, toxicity, or a multi-objective project profile. This is an
assay-local potency decision benchmark, not a broad biological-reasoning test.

\subsection{Action contract and constraints}

For task $c$, let $E_c$ be the visible support evidence, $P_c$ the candidate
pool, $V_c\subseteq P_c$ the candidates satisfying molecular eligibility, and
$k=10$. A system $m$ emits a schema-conforming action for task $c$ containing
an ordered sequence $S_m(c)$ of candidate IDs. The action is valid exactly when
the task identity matches, it has length $k$, it contains unique members of
$P_c$, it excludes support identities, and every selected candidate lies in
$V_c$.

Version 0.1.0 candidate eligibility requires an RDKit-parseable structure,
molecular weight at most 500 Da, and RDKit Crippen cLogP at most 4.5. The
500-Da limit follows the familiar rule-of-five upper bound
\cite{lipinski1997}. The prespecified 4.5 cLogP cutoff is a simplified benchmark
threshold, 0.5 below the conventional rule-of-five value of 5.0; it was fixed
before model evaluation. Across all 19,588
official candidate records, every structure parsed, 10,039 passed both limits,
2,433 exceeded molecular weight only, 3,765 exceeded cLogP only, and 3,351
exceeded both. The configured forbidden-substructure list was empty. These
thresholds provide transparent, reproducible, non-trivial constraints for
testing specification adherence. They are benchmark design choices, not
definitions of drug-likeness, oral bioavailability, synthesis feasibility,
medicinal-chemistry suitability, or safety.

\begin{figure}[H]
\centering
\fbox{\begin{minipage}{0.94\linewidth}
\centering
\textbf{Wider future loop (not evaluated)}\\[3pt]
design $\longrightarrow$ make $\longrightarrow$ test $\longrightarrow$ analyse
$\longrightarrow$ update\\[8pt]
\textbf{SpecGuard-Chem action boundary}\\[3pt]
\begin{tabular}{c@{\quad$\longrightarrow$\quad}c@{\quad$\longrightarrow$\quad}c}
support evidence + finite pool & filter, estimate, rank & ordered top-10 assay batch \\
constraints + budget & system under test & hidden retrospective outcomes
\end{tabular}
\end{minipage}}
\caption{The benchmark isolates one transition at which a model-generated
ranking becomes a finite experimental commitment. The surrounding autonomous
laboratory loop is motivation, not an evaluated claim.}
\label{fig:action_boundary}
\end{figure}

\section{Systems and interfaces}

\subsection{Non-language comparators}

We include a random-valid selector; a descriptor-desirability heuristic;
similarity to the most active support compound using Morgan fingerprints
\cite{rogers2010ecfp}; and random-forest, gradient-boosting, and linear
support-vector regressors trained separately for each card on only its 50
visible support measurements \cite{pedregosa2011sklearn,tropsha2010qsar}.
Structures are represented by 2,048-bit radius-2 Morgan fingerprints. Random
forest uses 100 trees, minimum leaf size 1, and seed 7. Gradient boosting uses
the locked scikit-learn defaults and seed 7. Linear SVR uses $C=1.0$, a linear
kernel, and sparse-compatible standard scaling without centering. No
card-specific hyperparameter optimization is performed. Each model predicts
all feasible candidates, which are sorted by predicted activity with candidate
ID as a deterministic tie-break. A non-deployable oracle ranks feasible
candidates using hidden outcomes and defines attainable headroom. Parsing,
descriptors, and fingerprints use RDKit \cite{rdkit2024}; exact software
versions are recorded in the supplement and lockfile.

\subsection{LLM interfaces and frozen conditions}

The final matrix crosses 91 cards, two interfaces, and three model conditions
(546 raw provider requests). Each request contains one complete card: task and
assay identifiers; the output contract; the 50 support IDs, structures, and
measured pChEMBL values; the ordered candidate IDs and structures; constraints;
and budget. The \emph{basic} interface also supplies candidate molecular weight
and cLogP. The \emph{descriptor-enriched} interface adds precomputed TPSA,
hydrogen-bond donor/acceptor counts, and rotatable-bond count. Hidden candidate
activities and external search are absent from both interfaces. Candidate
ordering and prompt wording are fixed within each condition.

The frozen conditions are OpenAI
\path{gpt-5.5-2026-04-23} with low reasoning effort, Anthropic
Claude Opus 4.8 without extended thinking, and
\path{deepseek-v4-pro} with thinking disabled, as selected and checked on 16
July 2026 \cite{openai2026gpt55,anthropic2026models,deepseek2026pricing}.
Every condition allows at most 4,096 output tokens. Temperature is not set and
therefore follows the provider default; no provider generation seed is set.
OpenAI and DeepSeek use provider JSON-object response modes. Anthropic uses an
explicit JSON-only prompt followed by deterministic extraction and validation
of one JSON object. SDK retries are disabled, and the matrix contains exactly
one successful recorded provider attempt per request. Thus ``raw'' below means
the parsed pre-repair output of that single response, not necessarily a
provider-native structured object.

Exact requests were frozen before execution. Configured and returned model
identifiers, response identifiers, settings, token counts, latency, finish
reasons, and raw content are retained. A conservative estimate placed the full
uncached matrix at \$106.06, with an explicit \$120 stop gate. The authorized
execution completed with \TotalLLMCost{} dollars of usage-derived token-pricing
cost. No call occurs unless the command also receives
\texttt{--allow-external}; exact replay uses the committed cache.

\subsection{Post-hoc deterministic repair}

A \emph{guarded LLM system} is the LLM plus deterministic software that checks
and, when needed, repairs its action. Here ``guarded'' refers only to compliance
with the benchmark contract; it does not imply biological, chemical, or
clinical safety. Raw and guarded behavior are evaluated from the same provider
response. The validator checks JSON shape, task identity, selection count,
candidate membership, duplicates, support exclusion, and candidate constraints;
it never reads activity outcomes. Valid, unique raw IDs are retained in
serialized response order, and unoccupied positions are filled from the fixed
rules-only ranking before revalidation. Schema or rank correction can therefore
count as a repaired card without replacing a candidate identity. The raw
response, raw issues, repaired output, and repair attribution remain stored.
This design avoids conflating a different repair prompt or an additional paid
model call with deterministic harness behavior. The outcome-blind fallback
formula is given in the supplement; it computes TPSA even for the basic
interface, so guarded results include software-only descriptor access that the
corresponding raw LLM did not receive.

\section{Evaluation}

\subsection{Utility, ranking, and regret}

Let $s_{mcr}$ be the candidate ID returned by system $m$ for card $c$ at
position $r$, with a missing symbol when that position was not supplied. Define
$a_{mcr}$ as $y_{s_{mcr},c}$ only when the position exists, its ID is a feasible
candidate with an available hidden outcome, and that ID has not occurred at an
earlier scored position; otherwise set $a_{mcr}=0$. Only the first $k$ returned
positions are scored. \emph{Feasible utility} is therefore
\begin{equation}
U_m(c)=\sum_{r=1}^{k}a_{mcr}.
\end{equation}
Thus the first occurrence of an otherwise valid ID can score, whereas later
duplicates, out-of-pool or infeasible IDs, missing outcomes, missing positions,
and positions beyond $k$ receive no credit. Zero is an executability penalty,
not an imputed activity. Because every card requires ten positions,
$U_m(c)/10$ is the mean pChEMBL recovered per required slot.

All source labels share the negative-log molar pChEMBL scale and
higher-is-better direction, but IC50, Ki, EC50, Potency, assay conditions, and
activity distributions are not biologically interchangeable. The primary
comparison is therefore paired within card. For systems $A$ and $B$, we first
calculate $\Delta_c=U_A(c)-U_B(c)$ on the same assay, support evidence, and
candidate pool, then equally weight the 91 card-level differences. For valid
fixed-size actions, an assay-specific additive offset cancels from this
contrast. Aggregate utility is a benchmark summary, not a clinical or
cross-assay physical quantity.

NDCG@$k$ is a secondary, within-card normalized measure of ordered ranking
quality relative to that card's ideal feasible top-10 ordering
\cite{jarvelin2002dcg}; 1.0 is ideal. Constrained regret is
\begin{equation}
R_m(c)=\sum_{i\in O_c}y_{ic}-U_m(c),
\end{equation}
where $O_c$ is the feasible hidden-outcome top-$k$. Lower regret is better.

\subsection{Action validity and operational outcomes}

Whole-action validity is a binary card-level measure: it equals one only when
the complete output has no contract issue, including schema shape, task
identity, exact size, in-pool membership, uniqueness, support exclusion, and
candidate eligibility. We separately report the valid-selection fraction
(\texttt{compliance\_rate}), defined as valid selected entries divided by $k$,
as a partial-credit diagnostic rather than a synonym for action validity. For
example, ten feasible selections returned under the wrong task identifier have
selection compliance one but action validity zero. For every LLM interface we
report raw metrics and the corresponding post-hoc-repaired metrics, plus repair
rate, fallback-supplied positions, and violation taxonomy. Latency, input/output
tokens, cache status, and estimated or actual cost are reported per condition.

\subsection{Bootstrap analysis}

Every resampling unit is a complete decision card. Marginal 95\% confidence
intervals for a system mean use 1,000 bootstrap samples from NumPy's default
random-number generator with seed 7. Each sample draws 91 cards with
replacement; its mean is calculated, and the 2.5th and 97.5th percentiles form
the interval. Headline system contrasts use a paired bootstrap: the 91 matched
card-level differences are calculated first, then resampled with replacement
2,000 times using seed 13. The observed mean difference and the corresponding
percentile interval are reported in the prespecified direction, including
best-QSAR minus best-guarded-LLM. Pairing preserves shared card difficulty;
bootstrapping estimates uncertainty but does not make assay endpoints
biologically equivalent or correct task selection, dependence beyond the card,
or model-selection multiplicity. Because the displayed best systems were
identified on this benchmark, their intervals are descriptive post-selection
intervals, not multiplicity-adjusted confirmatory tests or Bayesian posterior
probabilities.

\section{Results}

\subsection{The corrected benchmark separates molecular selectors}

Table~\ref{tab:baselines} establishes the benchmark's useful dynamic range
before provider spending. Eligibility filtering alone is not sufficient:
rules-only utility is \RulesUtility{} and random-valid utility is
\RandomUtility{}. Similarity rises to \SimilarityUtility{}, while all three
assay-local QSAR models reach approximately 74.75--74.97. The oracle remains at
\OracleUtility{}, leaving \BestQSARRegret{} utility points of mean headroom
above the best deployable baseline. All rows have perfect action validity, so
their separation is attributable to ranking rather than malformed actions.

\begin{table}[H]
\centering
\caption{Corrected deterministic results on all 91 cards. Utility intervals are
task-level 95\% bootstrap percentile intervals; NDCG@10 is a point estimate.
Oracle is not a deployable system.}
\label{tab:baselines}
\small
\begin{tabularx}{\linewidth}{@{}Yrrrr@{}}
\toprule
\textbf{System} & \textbf{Mean utility (95\% CI)} & \textbf{NDCG@10} & \textbf{Regret} & \textbf{Action validity} \\
\midrule
\DeterministicBaselineRows
\bottomrule
\end{tabularx}
\end{table}

\subsection{Do LLMs add decision value?}

\ifllmresultsavailable
The strongest final LLM pipeline was the OpenAI basic interface plus
deterministic post-hoc repair, with mean feasible utility \BestLLMUtility{} and
NDCG@10 \BestLLMNDCG{}. Linear-SVR QSAR remained higher by
\BestQSARMinusLLMUtility{} mean card-level utility points (paired 95\% CI
\BestQSARMinusLLMUtilityCILow{}--\BestQSARMinusLLMUtilityCIHigh{}). The
corresponding QSAR-minus-guarded-LLM NDCG@10 difference was 0.0087 (95\% CI
0.0010--0.0165). These two outcomes support a QSAR-family lead rather than added
decision value from the tested LLM pipelines.

The underlying pre-repair OpenAI response stream had mean feasible utility
\BestLLMRawUtility{} and issued a fully valid action on 72 of 91 cards
(\BestLLMActionValidity{}). Guarding raised utility to \BestLLMUtility{}, but
this is not the performance of the LLM alone. Repair was invoked on 19 cards;
only ten cards received a replacement identity, and deterministic fallback
supplied 14 of 910 final positions (1.54\%). Nine repairs corrected metadata,
schema, or ranking details without replacing an identity.

Harness contribution was much larger elsewhere
(Table~\ref{tab:repair_attribution}). Fallback supplied 25.16--27.69\% of
Anthropic final positions and 39.01--58.68\% of DeepSeek positions. Thirty
DeepSeek basic-interface cards received all ten final identities from fallback.
Those guarded scores primarily characterize a model-plus-deterministic-harness
system. Descriptor enrichment did not improve the best-performing OpenAI
condition, and effects varied by provider rather than supporting a general
descriptor benefit. The six raw conditions cost \TotalLLMCost{} dollars in
total under provider-reported usage and frozen token-pricing rules; the best
condition accounted for \BestLLMCost{} dollars. Repair added no provider calls.
\else
The canonical six-condition comparison failed the manuscript generator's
completeness or provenance gate. No partial, pilot-only, substituted, or
historical paper-50 number is reported in its place.
\fi

\begin{table}[H]
\centering
\caption{Pre-repair output and deterministic fallback contribution. A repaired
card can have zero replaced identities when only schema, task metadata, or rank
normalization changed. Final guarded rows are model-plus-harness results.}
\label{tab:repair_attribution}
\scriptsize
\begin{tabularx}{\linewidth}{@{}Yrrrr@{}}
\toprule
\textbf{Provider, interface} &
\textbf{Cards repaired} &
\textbf{Cards with identity replacement} &
\textbf{Fallback positions} &
\textbf{All 10 from fallback} \\
\midrule
\RepairAttributionCompactRows
\bottomrule
\end{tabularx}
\end{table}

\begin{figure}[H]
\centering
\includegraphics[width=0.92\linewidth]{../figures/v0.1.0/compliance_utility_frontier.png}
\caption{Action utility versus whole-action validity. Crosses are raw LLM
outputs; open triangles are deterministic post-hoc-repaired outputs from the
same responses. Connecting lines show harness intervention, not an additional
model call; repaired points must not be interpreted as model-only performance.}
\label{fig:utility_validity_frontier}
\end{figure}

\section{Discussion}

\subsection{A benchmark with a deliberately simple core}

The formal core of this task is filter--predict--rank--allocate. This simplicity
does not establish a new molecular-learning problem, and we do not claim
otherwise. It does create a precise evaluation boundary at the point
where a model's scores become an experimental commitment. That boundary makes
invalid actions, opportunity cost, deterministic repair, system inputs, and
operational cost measurable together.

The corrected evidence shows that the benchmark has independent legs:
constraints do not collapse the task, transparent molecular models improve
materially on heuristic selection, and oracle headroom remains. The tested LLMs
did not beat the best assay-local QSAR systems. For this bounded action, a
decision system should therefore retain the smaller specialized ranker as the
strong default and treat an LLM, if used, as an auditable component rather than
assuming it is the scientific decision engine.

\subsection{Relationship to existing evidence}

The strength of per-card QSAR is consistent with CARA's lead-optimization
analysis, where models trained on assay-relevant examples performed strongly
and performance varied among assays \cite{tian2024cara}. Variation among
providers and interfaces is likewise consistent with chemistry benchmarks
showing that LLM performance depends on task and evaluation setup
\cite{guo2023chemllmbench,mirza2025chembench}. Mol-Instructions, LLM4SD, and
ChemCrow demonstrate useful roles for domain instruction data, LLM-derived
features, and tool coordination
\cite{fang2024molinstructions,zheng2025llm4sd,bran2024chemcrow}; none implies
that a directly prompted LLM should outperform a regressor fitted to local
assay measurements. Our result is narrower: on this corrected retrospective
top-10 benchmark, guarded pipelines produced executable shortlists, but the
strongest conventional QSAR retained a paired advantage.

\subsection{What the benchmark does not establish}

This is a retrospective computational benchmark, not prospective lead
optimization. It includes every one of the 91 \texttt{LO\_All} tasks that can
support the declared top-10 action, but it does not represent CARA
virtual-screening tasks, the nine tasks with fewer than ten feasible candidates,
or the distribution of industrial projects. CARA also inherits ChEMBL
measurement and curation heterogeneity \cite{zdrazil2024chembl}. Bootstrap
intervals describe variation among these observed cards and do not license
generalization beyond them.

The objective is assay-local potency, not selectivity, mechanism, uncertainty,
ADMET, toxicity, pharmacokinetics, synthesis, novelty, availability, or
therapeutic value. The candidate pool is fixed and no structure is generated.
The molecular thresholds are simplified eligibility rules; no structural alert
was active, and no medicinal chemist comparator was included. Consequently,
``feasible'' means feasible under the benchmark's rules, not suitable for
synthesis or development.

Each provider condition retains one response per card, so repeated-call
variation is not estimated. Results are specific to the dated models, prompts,
interfaces, parsing, settings, and deterministic fallback. CARA and ChEMBL are
public: candidate outcomes were withheld at prompt time and from non-oracle
ranking code, but prior model exposure to public source data cannot be excluded.
One-shot allocation also does not test learning across rounds. These limits
preclude claims about prospective laboratory performance, autonomous-laboratory
readiness, or LLMs in drug discovery generally.

\section{Reproducibility and release}

The frozen working release separates public cards from scorer outcomes and
includes source and transform provenance, exclusions, exact requests, model
configurations, traces, per-card scores, aggregate tables, costs, environment
locks, licenses, citation metadata, and reproduction commands. Historical
invalid results remain immutable but are explicitly superseded. A semantic
release tag and the final top-level manifest/checksum freeze remain separate
release-administration actions.

\section{Conclusion}

Across all 91 eligible CARA lead-optimization cards, the strongest
assay-specific QSAR model outperformed the strongest tested guarded LLM
pipeline on the paired primary outcome. Deterministic validation made final
actions executable, but fallback supplied between 1.54\% and 58.68\% of final
shortlist positions across conditions. Guarded performance must therefore be
reported as model-plus-harness performance alongside the pre-repair model
result. This study supports strong conventional comparators, paired card-level
analysis, and explicit repair attribution when evaluating constrained LLM
decision systems; it does not establish superiority or general utility of LLMs
for prospective drug discovery.

\section*{Data, code, and AI-use statement}

Release artifacts and exact reproduction commands accompany the repository.
Code is MIT licensed; CARA and CARA-derived data retain CC BY 4.0
terms. ChatGPT/Codex assisted with project framing, code and artifact auditing,
manuscript structuring/editing, and LaTeX compilation. All numerical claims are
checked against local machine-readable artifacts. No patient or clinical data
are used. Author list, contributions, funding, competing interests, and
venue-specific disclosure wording require human confirmation before submission.

\begingroup
\small
\bibliographystyle{plain}
\bibliography{references}
\endgroup

\end{document}
